\documentclass[a4paper,12pt]{article}

\usepackage[T2A]{fontenc}
\usepackage[utf8]{inputenc}
\usepackage[russian]{babel}
\usepackage{amsmath,amssymb,amsthm,mathtools}
\usepackage{helvet}
\renewcommand{\familydefault}{\sfdefault}
\usepackage{icomma}
\usepackage[a4paper,margin=2.2cm,headheight=25pt]{geometry}
\usepackage{microtype}
\usepackage{tikz}
\usetikzlibrary{arrows.meta,positioning,calc}
\usepackage{tabularx,booktabs}
\usepackage{enumitem}
\usepackage{xcolor}
\usepackage{hyperref}
\usepackage{parskip}
\usepackage{fancyhdr}
\usepackage{setspace}
\usepackage{titlesec}

\definecolor{accent}{HTML}{008080}
\definecolor{darkaccent}{HTML}{004D4D}
\definecolor{textgray}{HTML}{333333}
\definecolor{softgray}{HTML}{EFF4F4}
\definecolor{linegray}{HTML}{AAB8B8}

\hypersetup{
  colorlinks=true,
  linkcolor=darkaccent,
  urlcolor=accent,
  citecolor=darkaccent
}

\onehalfspacing
\setlength{\parindent}{0pt}
\setlength{\tabcolsep}{5pt}
\renewcommand{\arraystretch}{1.35}
\emergencystretch=2em
\raggedbottom

\titleformat{\section}
  {\Large\bfseries\color{darkaccent}}
  {\thesection.}{0.55em}{}
\titleformat{\subsection}
  {\large\bfseries\color{accent}}
  {\thesubsection.}{0.5em}{}
\titlespacing*{\section}{0pt}{1.4em}{0.65em}
\titlespacing*{\subsection}{0pt}{1.1em}{0.45em}

\pagestyle{fancy}
\fancyhf{}
\fancyhead[L]{\small Дифференцированная схема кредитования}
\fancyhead[R]{\small Николь Саркисьянц}
\fancyfoot[C]{\small\thepage}
\renewcommand{\headrulewidth}{0.4pt}
\renewcommand{\headrule}{%
  \hbox to\headwidth{%
    \color{linegray}\leaders\hrule height \headrulewidth\hfill
  }%
}

\newcolumntype{Y}{>{\centering\arraybackslash}X}
\newcommand{\rub}{\,\text{руб.}}
\newcommand{\mln}{\,\text{млн руб.}}
\newcommand{\thousand}{\,\text{тыс. руб.}}
\newcommand{\Total}{A_{\text{общ}}}
\newcommand{\Over}{A_{\text{пер}}}

\begin{document}

\begin{titlepage}
  \thispagestyle{empty}
  \centering
  \vspace*{2.6cm}

  {\Large\bfseries\color{accent} ЧЕК-ЛИСТ\par}
  \vspace{0.8cm}

  {\Huge\bfseries\color{darkaccent}
  Дифференцированная схема кредитования\par}

  \vspace{0.55cm}

  {\Large
  Кредитные задачи ЕГЭ: таблица, формулы, прогрессия, переплата\par}

  \vspace{2.2cm}

  {\large Лёвин Артём Александрович\par}
  \vspace{0.25cm}
  {\large эксклюзивно для Николь Саркисьянц\par}

  \vfill

  {\large \today\par}
  \vspace{1.3cm}

  \begin{tikzpicture}
    \draw[accent!58,line width=1.15pt]
      (0,0)
      .. controls (2.2,0.55) and (4.1,-0.25) .. (6.2,0.15)
      .. controls (8.2,0.52) and (10.5,-0.18) .. (12.4,0.18)
      .. controls (13.7,0.42) and (14.6,0.25) .. (15.2,0.05);
  \end{tikzpicture}

  \vspace*{0.7cm}
\end{titlepage}

\section{Суть дифференцированной схемы}

В дифференцированной схеме тело кредита погашается равными частями.
Проценты каждого периода начисляются на остаток долга к началу этого
периода. Остаток постепенно уменьшается, поэтому платежи образуют
убывающую арифметическую прогрессию.

\subsection{Обозначения}

\begin{table}[ht]
\centering
\caption{Основные величины}
\begin{tabularx}{\linewidth}{@{}lX@{}}
\toprule
Обозначение & Смысл \\
\midrule
$S$ & первоначальная сумма кредита, $S>0$; \\
$n$ & количество платежных периодов, $n\in\mathbb{N}$; \\
$p$ & процентная ставка за один период, записанная десятичной дробью; \\
$r$ & та же ставка в процентах, поэтому
      $p=\dfrac{r}{100}$; \\
$P_k$ & платеж в конце $k$-го периода,
        $k=1,2,\ldots,n$; \\
$B_k$ & остаток долга после $k$-го платежа;
        $B_0=S$, $B_n=0$. \\
\bottomrule
\end{tabularx}
\end{table}

За каждый период погашается одна и та же часть тела кредита:
\[
\boxed{\frac{S}{n}}.
\]

\subsection{Порядок операций в каждом периоде}

\begin{center}
\begin{tikzpicture}[
  node distance=9mm and 7mm,
  every node/.style={font=\small,align=center},
  stage/.style={
    draw=accent!65,
    rounded corners=2pt,
    inner sep=6pt,
    minimum height=12mm
  },
  arr/.style={
    -{Latex[length=2.5mm]},
    line width=0.8pt,
    color=darkaccent
  }
]
\node[stage] (b0) {остаток долга\\в начале периода};
\node[stage,right=of b0] (int) {начисление\\процентов};
\node[stage,right=of int] (debt) {общая\\задолженность};
\node[stage,right=of debt] (pay) {платеж};
\node[stage,right=of pay] (b1) {новый остаток\\долга};

\draw[arr] (b0) -- (int);
\draw[arr] (int) -- (debt);
\draw[arr] (debt) -- (pay);
\draw[arr] (pay) -- (b1);
\end{tikzpicture}
\end{center}

Очередность особенно важна в задачах с датами: сначала проценты
увеличивают долг, затем заемщик вносит платеж.

\section{Абстрактная таблица кредита}

К началу $k$-го периода уже погашено $k-1$ равных частей тела кредита.
Поэтому начальный долг этого периода равен
\[
S-\frac{(k-1)S}{n}
=
\frac{S(n-k+1)}{n}.
\]

\begin{table}[ht]
\centering
\caption{Математическая модель дифференцированной схемы}
\small
\begin{tabularx}{\linewidth}{@{}cYYYY@{}}
\toprule
Период
&
Долг в начале
&
Начисленные проценты
&
Платеж
&
Остаток после платежа
\\
\midrule
$1$
&
$S$
&
$Sp$
&
$Sp+\dfrac{S}{n}$
&
$\dfrac{S(n-1)}{n}$
\\[2mm]

$2$
&
$\dfrac{S(n-1)}{n}$
&
$\dfrac{S(n-1)}{n}p$
&
$\dfrac{S(n-1)}{n}p+\dfrac{S}{n}$
&
$\dfrac{S(n-2)}{n}$
\\[2mm]

$k$
&
$\dfrac{S(n-k+1)}{n}$
&
$\dfrac{S(n-k+1)}{n}p$
&
$\dfrac{S(n-k+1)}{n}p+\dfrac{S}{n}$
&
$\dfrac{S(n-k)}{n}$
\\[2mm]

$n$
&
$\dfrac{S}{n}$
&
$\dfrac{Sp}{n}$
&
$\dfrac{Sp}{n}+\dfrac{S}{n}$
&
$0$
\\
\bottomrule
\end{tabularx}
\end{table}

\subsection{Формулы остатка}

После $k$-го платежа остается
\[
\boxed{B_k=\frac{S(n-k)}{n}},
\qquad
k=0,1,\ldots,n.
\]

К началу $k$-го периода остаток равен $B_{k-1}$:
\[
\boxed{B_{k-1}=\frac{S(n-k+1)}{n}}.
\]

Полезная проверка крайних случаев:
\[
B_0=S,
\qquad
B_n=0.
\]

\section{Платежи и арифметическая прогрессия}

\subsection{Платеж в произвольном периоде}

Платеж складывается из начисленных процентов и постоянной части тела
кредита:
\[
\boxed{
P_k=
\frac{S(n-k+1)}{n}p+\frac{S}{n}
}.
\]

Равносильная форма:
\[
P_k=
\frac{S}{n}\bigl(1+p(n-k+1)\bigr).
\]

Первый и последний платежи:
\[
\boxed{P_1=Sp+\frac{S}{n}},
\qquad
\boxed{P_n=\frac{S}{n}(1+p)}.
\]

Последний платеж является наименьшим при $p>0$.

\subsection{Разность соседних платежей}

\begin{align*}
P_{k+1}-P_k
&=
\left(
\frac{S(n-k)}{n}p+\frac{S}{n}
\right)
-
\left(
\frac{S(n-k+1)}{n}p+\frac{S}{n}
\right)
\\
&=
-\frac{Sp}{n}.
\end{align*}

Следовательно, платежи образуют убывающую арифметическую прогрессию
с постоянной разностью
\[
\boxed{d=-\frac{Sp}{n}}.
\]

Каждый следующий платеж меньше предыдущего на
\[
\boxed{\frac{Sp}{n}}.
\]

\section{Общая сумма выплат и переплата}

\subsection{Вывод через сумму процентов}

Сумма всех погашенных частей тела кредита равна $S$:
\[
\underbrace{
\frac{S}{n}+\frac{S}{n}+\cdots+\frac{S}{n}
}_{n\text{ раз}}
=
S.
\]

Сумма начисленных процентов:
\begin{align*}
Sp
+\frac{S(n-1)}{n}p
+\frac{S(n-2)}{n}p
+\cdots
+\frac{Sp}{n}
&=
\frac{Sp}{n}
\bigl(n+(n-1)+\cdots+1\bigr)
\\
&=
\frac{Sp}{n}\cdot\frac{n(n+1)}{2}
\\
&=
\frac{Sp(n+1)}{2}.
\end{align*}

Отсюда получаем ключевые формулы.

\subsection{Основные формулы}

Общая сумма выплат банку:
\[
\boxed{
\Total=S+\frac{Sp(n+1)}{2}
}
\]
или
\[
\boxed{
\Total=
S\left(1+\frac{p(n+1)}{2}\right)
}.
\]

Переплата:
\[
\boxed{
\Over=
\Total-S=
\frac{Sp(n+1)}{2}
}.
\]

Процент переплаты относительно первоначальной суммы кредита:
\[
\boxed{
q=
\frac{\Over}{S}\cdot100\%
=
\frac{p(n+1)}{2}\cdot100\%
}.
\]

При ставке $r\%$ за период формула упрощается:
\[
\boxed{
q=\frac{r(n+1)}{2}\%
}.
\]

Первоначальная сумма $S$ сокращается, поэтому процент переплаты зависит
от ставки и количества периодов.

\subsection{Второй способ: сумма арифметической прогрессии}

Поскольку $P_1,P_2,\ldots,P_n$ образуют арифметическую прогрессию,
\[
\Total=
\frac{P_1+P_n}{2}\cdot n.
\]

Эта формула особенно полезна, когда известны первый платеж, последний
платеж и общая сумма выплат.

\section{Как распознать схему по условию}

Признаки дифференцированной схемы:
\begin{itemize}[leftmargin=1.5em,itemsep=0.25em]
  \item долг уменьшается каждый период на одну и ту же величину;
  \item основная сумма долга погашается равными долями;
  \item платежи постепенно уменьшаются;
  \item проценты начисляются на остаток долга;
  \item наименьший платеж приходится на последний период.
\end{itemize}

Формулировка «выплаты осуществляются равными платежами» указывает
на аннуитетную схему.

\section{Алгоритм решения кредитной задачи}

\begin{enumerate}[
  leftmargin=1.8em,
  label=\arabic*.,
  itemsep=0.35em
]
  \item Определите единицу периода: месяц, год или другой указанный
        промежуток.

  \item Введите обозначения $S$, $n$, $p$ и при необходимости
        $P_k$, $B_k$.

  \item Переведите процентную ставку в десятичную дробь:
        $p=r/100$.

  \item Зафиксируйте равную часть тела кредита $S/n$.

  \item Запишите минимум строки для первого, второго и последнего
        периодов либо строку для произвольного $k$.

  \item Выразите требуемую величину: остаток, платеж, общую сумму,
        переплату или процент переплаты.

  \item Используйте арифметическую прогрессию, когда требуется сумма
        платежей или сумма остатков.

  \item Решите полученное уравнение и проверьте ограничения:
        $S>0$, $n\in\mathbb{N}$, $p\ge0$.

  \item Проверьте размерность ответа и смысл результата.
\end{enumerate}

\section{Минимальное оформление развернутого решения}

Для полноценного обоснования на экзамене достаточно компактной
математической модели:
\begin{enumerate}[
  leftmargin=1.8em,
  label=\arabic*.,
  itemsep=0.3em
]
  \item обозначения $S$, $n$, $p$;
  \item строки таблицы для периодов $1$, $2$, $n$ либо общая строка
        для $k$;
  \item формула платежа $P_k$;
  \item указание на арифметическую прогрессию;
  \item вывод нужной формулы;
  \item числовая подстановка и ответ.
\end{enumerate}

Готовую формулу следует сопровождать выводом из модели кредита.
Это делает решение проверяемым и логически завершенным.

\section{Типичные ошибки}

\begin{itemize}[leftmargin=1.5em,itemsep=0.3em]
  \item Начисление процентов на первоначальную сумму $S$ во всех
        периодах. Основа начисления меняется вместе с остатком долга.

  \item Использование остатка после $k$-го платежа при расчете
        процентов за $k$-й период. Нужен остаток к началу периода
        $B_{k-1}$.

  \item Подстановка $r$ вместо $p$. При $r=12\%$ в формулах
        используется $p=0{,}12$.

  \item Смешение общей суммы выплат и переплаты:
        $\Over=\Total-S$.

  \item Потеря множителя $n$ в формуле суммы арифметической
        прогрессии.

  \item Ошибка в числе слагаемых: последовательность
        $1+2+\cdots+n$ содержит ровно $n$ членов.

  \item Неверный индекс остатка: после $k$ платежей остается
        $S(n-k)/n$.

  \item Получение дробного значения $n$ без анализа. Количество
        платежных периодов должно быть натуральным числом.

  \item Отсутствие проверки последнего остатка. После $n$-го платежа
        долг обязан стать равным нулю.
\end{itemize}

\section{Короткие примеры}

\subsection*{Пример 1. Процент переплаты}

Кредит погашается за $9$ месяцев по ставке $12\%$ в месяц. Тогда
\[
q=
\frac{12(9+1)}{2}\%
=
60\%.
\]

\subsection*{Пример 2. Последний платеж}

При $S=3\mln$, $p=0{,}2$ и $n=15$:
\[
P_{15}
=
\frac{3}{15}(1+0{,}2)
=
0{,}24\mln.
\]

\subsection*{Пример 3. Остаток после нескольких платежей}

При $S=1{,}2\mln$, $n=12$ после пятого платежа:
\[
B_5
=
\frac{1{,}2(12-5)}{12}
=
0{,}7\mln.
\]

\clearpage

\section*{Тренировочный блок}

Во всех задачах проценты начисляются на остаток долга, после чего
вносится платеж. Ответы запишите с указанными единицами измерения.

\subsection*{Средний уровень}

\begin{enumerate}[
  leftmargin=1.8em,
  label=\textbf{\arabic*.},
  itemsep=0.9em
]
  \item Кредит $900\thousand$ погашается дифференцированными платежами
        за $9$ месяцев. Ставка составляет $2\%$ в месяц. Найдите
        постоянную часть тела кредита, первый платеж, последний платеж,
        общую сумму выплат и процент переплаты.

  \item Кредит $1{,}2\mln$ выдан на $12$ месяцев под $1\%$ в месяц.
        Найдите пятый платеж и остаток долга после пятого платежа.

  \item Сергей взял кредит на $9$ месяцев. В конце каждого месяца долг
        сначала увеличивается на $12\%$, затем уменьшается на внесенный
        платеж. Остаток долга каждый месяц уменьшается на одну и ту же
        величину. Сколько процентов от суммы кредита составит
        переплата?
\end{enumerate}

\subsection*{Уровень выше среднего}

\begin{enumerate}[
  leftmargin=1.8em,
  label=\textbf{\arabic*.},
  start=4,
  itemsep=0.9em
]
  \item Кредит $3\mln$ выдан на $n$ лет под $20\%$ годовых.
        Погашение происходит по дифференцированной схеме. Наименьший
        платеж равен $240\thousand$. Найдите $n$, общую сумму выплат
        и переплату.

  \item Кредит погашается за $10$ лет под $20\%$ годовых.
        Последний платеж равен $180\thousand$. Найдите первоначальную
        сумму кредита, первый платеж и общую сумму выплат.

  \item Кредит $1{,}6\mln$ погашается за $5$ лет. Общая сумма выплат
        составила $2{,}08\mln$. Найдите годовую процентную ставку
        и третий платеж.
\end{enumerate}

\clearpage

\subsection*{Уровень выше среднего --- продолжение}

\begin{enumerate}[
  leftmargin=1.8em,
  label=\textbf{\arabic*.},
  start=7,
  itemsep=0.9em
]
  \item Кредит погашается за $24$ месяца. Переплата составляет
        $18{,}75\%$ от первоначальной суммы. Найдите месячную
        процентную ставку.

  \item Кредит $2{,}4\mln$ погашается за $12$ месяцев под $1{,}5\%$
        в месяц. На сколько рублей уменьшается каждый следующий платеж?
        Найдите сумму платежей с четвертого по восьмой месяц
        включительно.

  \item Кредит погашается за $10$ лет под $10\%$ годовых. Общая сумма
        выплат равна $1{,}55\mln$. Найдите первоначальную сумму кредита
        и седьмой платеж.
\end{enumerate}

\subsection*{Сложный уровень}

\begin{enumerate}[
  leftmargin=1.8em,
  label=\textbf{\arabic*.},
  start=10,
  itemsep=0.9em
]
  \item В дифференцированной схеме первый платеж равен
        $420\thousand$, последний платеж равен $150\thousand$,
        а общая сумма выплат равна $2{,}85\mln$. Найдите количество
        платежей, первоначальную сумму кредита, процентную ставку
        за период и процент переплаты.
\end{enumerate}

\clearpage

\section*{Ответы к тренировочному блоку}

\begin{table}[ht]
\centering
\caption{Краткие ответы}
\begin{tabularx}{\linewidth}{@{}cX@{}}
\toprule
№ & Ответ \\
\midrule
1 &
$100\thousand$;
$118\thousand$;
$102\thousand$;
$990\thousand$;
$10\%$.
\\

2 &
$108\thousand$;
$700\thousand$.
\\

3 &
$60\%$.
\\

4 &
$n=15$;
$7{,}8\mln$;
$4{,}8\mln$.
\\

5 &
$1{,}5\mln$;
$450\thousand$;
$3{,}15\mln$.
\\

6 &
$10\%$ годовых;
$416\thousand$.
\\

7 &
$1{,}5\%$ в месяц.
\\

8 &
На $3\thousand$;
$1{,}105\mln$.
\\

9 &
$1\mln$;
$140\thousand$.
\\

10 &
$n=10$;
$S=1{,}2\mln$;
$25\%$ за период;
$137{,}5\%$.
\\
\bottomrule
\end{tabularx}
\end{table}

\end{document}